\section{Chứng minh phi tương tác và phép biến đổi Fiat–Shamir}

(Non-Interactive Proofs and Fiat-Shamir)

\subsection{(1) [20 điểm] Phép biến đổi Fiat–Shamir (Fiat–Shamir Transform)}

\subsubsection{(a) [10 điểm] Từ tương tác sang phi tương tác (From Interactive to Non-Interactive)}

Bạn đang chuyển đổi giao thức nhận dạng Schnorr thành một sơ đồ chữ ký số bằng phép biến đổi Fiat–Shamir.

\textbf{i. Trình bày toàn bộ thuật toán ký cho thông điệp:}

``Transfer \$1000 to Bob''

Bao gồm cách tính thách thức $c = H(A, Y, m)$.

\textbf{ii. Giải thích vì sao việc đưa thông điệp $m$ vào trong hàm băm là cực kỳ quan trọng.}

Nếu thay vào đó sử dụng $c = H(A, Y)$, thì loại tấn công nào sẽ trở nên khả thi?

\vspace{0.5cm}

\subsubsection{(b) [10 điểm] Bảo mật trong mô hình Random Oracle}

Phép biến đổi Fiat–Shamir giả định rằng hàm băm hoạt động như một Random Oracle.

\textbf{i. Giải thích những thuộc tính của Random Oracle mà một hàm băm thực tế (như SHA-256) có thể không đạt được hoàn hảo.}

\textbf{ii. Mô tả ở mức cao Bổ đề Forking (Forking Lemma).}

Giải thích cách bổ đề này cho phép rút trích nhân chứng (witness) từ một kẻ tấn công có khả năng giả mạo chứng minh Fiat–Shamir.

\vspace{0.5cm}

\subsection{(2) [15 điểm] Mạch và Zero-Knowledge tổng quát (Circuits and General-Purpose ZK)}

Bạn muốn chứng minh rằng mình biết một giá trị $x$ sao cho
$\text{SHA-256}(x) = y$
với một $y$ đã cho.

\subsubsection{(a) [7.5 điểm] Giải thích vì sao việc tạo một giao thức Sigma riêng cho phát biểu này là rất khó khăn.}

Làm thế nào việc biểu diễn SHA-256 dưới dạng mạch (circuit) lại khiến chứng minh này trở nên khả thi?

\subsubsection{(b) [7.5 điểm] Ứng dụng của bạn cần chứng minh mệnh đề:}

``Tôi biết $x$ sao cho $x^3 + x + 5 = 35$ trong $\mathbb{Z}_{101}$''

Bạn sẽ chọn giao thức tùy chỉnh hay dạng mạch tổng quát?
Hãy giải thích lựa chọn của bạn.
